\FloatBarrier
\section{Development study of rollout objective and context dependence}
\label{sec:rollout-revision}
The completed dynamics-residual development experiment does not improve its
frozen Framewise donor: eligible success is 1/31 for both models on PushT and
8/30 versus 11/30 on Reacher. We therefore record a new, post-hoc development
protocol rather than promote that revision to the main test. All eight arms
have completed training and development planning in separate environment
campaigns. A $2\times2$
design crosses teacher-forced versus recursive prediction with inferred
versus constant-input dynamics context. Each of the four variants per
environment trains for 30 epochs at seed 0, freezing the same completed
Framewise calibrator, goal pathway, action encoder and predictor. Both
objectives supervise target frames 3--7 (indices \texttt{3:8}), including
the first prediction after the three-frame support. Teacher forcing uses
available observed history for each target; recursive training feeds back
predicted latents after support. Context uses support only. The constant-input
control retains the identical GRU/FiLM architecture but receives all-zero
support and executed-action inputs: its learned response can be nonzero,
and matching nominal parameter counts does not establish equal effective
capacity. All variants reuse the existing train/validation data and optimizer
schedule; validation-selected checkpoints minimize the same FP32 recursive
MSE averaged over all five query frames. Development planning uses the
existing 32 matched tasks per environment, 300 CEM candidates, 30 iterations,
30 elites and the fixed 50-native-step budget including support. All eight
prespecified arms are retained, and only validation error selects
checkpoints within this follow-up protocol.
Table~\ref{tab:rollout-revision-plan} reports source-validated status and
distinguishes incomplete training, trained checkpoints and completed planning;
missing measurements are
never substituted with zeros or partial planning outcomes.
The completed teacher-forced arms give mixed evidence about context dependence.
Inferred context ties the frozen donor on PushT (1/31 eligible successes) and
Reacher (11/30). Constant input scores 0/31 on PushT and 15/30 on Reacher,
compared with donor scores of 1/31 and 11/30, respectively. The latter is a
favorable development point estimate for the constant-input control, not
evidence that inferred context is necessary. The completed recursive/inferred
arms tie the donor on PushT (1/31) and score 13/30 versus 11/30 on Reacher
(a $+6.67$ percentage-point development difference). Recursive constant
context scores 0/31 on PushT and 13/30 on Reacher. Thus, inferred context
exceeds its constant-input control on PushT but ties the donor; on Reacher,
recursive inferred and constant contexts tie. Neither objective satisfies
the registered screening requirement to exceed both its constant-input
control and donor in both environments. Lower recursive validation error
does not establish better control, and these results do not establish a
benefit from episode-dependent dynamics inference. Repeated training seeds
and an untouched evaluation remain necessary to assess any subsequent method.
A revised context method is not promoted without evidence of learning beyond
its matched constant-input control; multiple training seeds and a fresh final
holdout unused during revision development are required before a broader
performance claim. These experiments remain separate from the original
three-seed campaign and its negative outcomes.

\IfFileExists{generated/rollout_revision.tex}{
% Generated only from source-validated rollout study records.
\begin{table}[H]\centering\small\setlength{\tabcolsep}{4pt}
\begin{tabular}{lllrrrr}\toprule
Environment & Objective & Context & Status & Val. MSE & Eligible & $\Delta$ (pp) \\\midrule
PushT & Teacher-forced & Inferred & Complete & 0.28759 & 1/31 & +0.00 \\
PushT & Teacher-forced & Constant & Complete & 0.28791 & 0/31 & -3.23 \\
PushT & Recursive & Inferred & Complete & 0.26388 & 1/31 & +0.00 \\
PushT & Recursive & Constant & Complete & 0.26876 & 0/31 & -3.23 \\
Reacher & Teacher-forced & Inferred & Complete & 0.00566 & 11/30 & +0.00 \\
Reacher & Teacher-forced & Constant & Complete & 0.00565 & 15/30 & \positivegain{+13.33} \\
Reacher & Recursive & Inferred & Complete & 0.00521 & 13/30 & \positivegain{+6.67} \\
Reacher & Recursive & Constant & Complete & 0.00525 & 13/30 & \positivegain{+6.67} \\
\bottomrule\end{tabular}
\caption{\textbf{Development objective/context study: protocol and completion status.} Eight seed-zero arms train for 30 extra epochs with a frozen Framewise donor. Planned means completed training is unavailable; trained means validated training is complete but planning is incomplete; complete requires all 32 matched development tasks. Validation MSE is the completed-history minimum of the common FP32 recursive criterion over targets 3--7, not a test score. Eligible success excludes support-only successes; $\Delta$ compares the same eligible tasks against the frozen Framewise donor. Bold green marks strictly positive observed differences, not statistical significance. Missing values are unavailable, never zero. Constant-input contexts can learn nonzero responses; nominal parameter matching does not establish equal effective capacity. Both objectives use the first target after support. Planning uses 300 CEM candidates, 30 iterations, 30 elites and the fixed native-action budget. This post-hoc development study does not establish a main-test, multi-seed, efficiency or context-necessity claim.}
\label{tab:rollout-revision-plan}\end{table}

}{% Planned experiment rows only. No measurements exist in this template.
\begin{table}[H]
\centering\small
\begin{tabular}{lllrr}
\toprule
Environment & Objective & Context & Eligible success & $\Delta$ vs donor (pp) \\
\midrule
PushT & Teacher-forced & Inferred & \missing & \missing \\
PushT & Teacher-forced & Constant & \missing & \missing \\
PushT & Recursive & Inferred & \missing & \missing \\
PushT & Recursive & Constant & \missing & \missing \\
\midrule
Reacher & Teacher-forced & Inferred & \missing & \missing \\
Reacher & Teacher-forced & Constant & \missing & \missing \\
Reacher & Recursive & Inferred & \missing & \missing \\
Reacher & Recursive & Constant & \missing & \missing \\
\bottomrule
\end{tabular}
\caption{\textbf{Planned development experiment; all results pending.}
Each row specifies one seed-zero, 30-epoch run with the same frozen Framewise
donor. Eligible success excludes tasks already solved during support;
$\Delta$ is the eventual eligible-success difference in percentage points
against that environment's unchanged donor (PushT: 1/31; Reacher: 11/30).
All dashes mean unavailable measurements, not zero performance. Teacher-forced
and recursive objectives share target frames 3--7 and a common recursive
validation criterion. Constant denotes a learned context response to zero
support/action inputs, not a zero residual. Beating the frozen donor alone
would not establish context necessity; the matched constant-input rows are
required controls.}
\label{tab:rollout-revision-plan}
\end{table}
}
